% !TeX document-id = {e2312f8c-b2d9-4f5f-a6e9-54d5ee345192}
% !TeX program = lualatex
% !BIB program = biber
% Lualatex is important to render TTF fonts; with pdflatex it's just the regular one
% ratio 16:9 -- https://tex.stackexchange.com/questions/14336/

% compile two versions, inspired by https://tex.stackexchange.com/a/1501
% use the script "compile-pdf.sh"
\newif\ifhandout
% if flags.tex does not exist, create an empty file to be able to compile in TeXstudio
\input{flags}

\ifhandout
\documentclass[12pt,aspectratio=169,handout]{beamer}
\else
\documentclass[12pt,aspectratio=169]{beamer}
\fi


% TODO change "leftfootertext" to your liking
\newcommand{\leftfootertext}{\insertsubtitle}  % just the \title{} text by default
%\newcommand{\leftfootertext}{RNNs and encoder-decoder architectures}  % Your name, for instance


% ------- RUB specifics ----------
% adjust for 16:9
% https://tex.stackexchange.com/questions/354022/modifying-the-margins-of-all-slides-in-beamer
\setbeamersize{text margin left=0.3cm,text margin right=4.5cm} 


% use Metropolis as the basis theme
\usetheme[subsectionpage=progressbar]{metropolis}
% blocks with background globally
\metroset{block=fill}


\usepackage{fontspec}
% RUB fonts need to be installed
% 'UprightFont = * Light' makes sure that the base font is RubFlama Light, which looks
% lighter than RubFlama Regular (would be too thick for slides)
\setsansfont[Scale=MatchLowercase, UprightFont = * Light, BoldFont = * Bold]{RubFlama}
%\setsansfont{Arial} % Open source alternative if you don't have RubFlama

% RUB color scheme
% Dark blue: 0; 53; 96; #003560
\definecolor{RUBDarkBlue}{RGB}{0, 53, 96}

% Light yellow (table fill, etc.); 238; 250; 196; #EEFAC4
\definecolor{RUBLightYellow}{RGB}{238, 250, 196}

%Light green: 141; 174; 16
\definecolor{RUBLightGreen}{RGB}{141, 174, 16}


\setbeamercolor{titlelike}{fg=RUBDarkBlue}
\setbeamercolor{subtitle}{fg=RUBLightGreen}
\setbeamercolor{separation line}{fg=RUBLightGreen}
\setbeamercolor{frametitle}{bg=white, fg=RUBDarkBlue}

% horizontal line on title page and sections
\setbeamercolor{alerted text}{fg=RUBLightGreen}


% Adjust footer bottom (too large by default)
\setbeamertemplate{footline}{%
  \begin{beamercolorbox}[wd=\textwidth, sep=2ex]{footline}%
    \usebeamerfont{page number in head/foot}%
    \usebeamertemplate*{frame footer}
    \hfill%
    \usebeamertemplate*{frame numbering}
  \end{beamercolorbox}%
}


% Lab name, numbering, etc. in footer
\setbeamertemplate{frame numbering}{TrustHLT --- Prof.\ Dr.\ Ivan Habernal \hspace*{1ex} \includegraphics[width=7em]{img/rub-logo.pdf}\hspace*{1ex}}

\setbeamertemplate{frame footer}{\hspace*{1ex}\insertframenumber \hspace*{2ex} \leftfootertext}

% adjust the background to be completely white
\setbeamercolor{background canvas}{bg=white}

% logos on the title page
\titlegraphic{%
	\begin{picture}(0,0)
		\put(435,0){\makebox(0,0)[rt]{\includegraphics[width=7em]{img/rub-logo.pdf}}}
		\put(435,-170){\makebox(0,0)[rt]{\includegraphics[width=4em]{img/logo-trusthlt.pdf}}}
		\put(435,-196){\makebox(0,0)[rt]{\includegraphics[width=9em]{img/logo-rctrust.pdf}}}
	\end{picture}%
}


% show TOC at every section start
\AtBeginSection{
	\frame{
		\vspace{2em}
		\sectionpage
		\hspace*{2.2em}\begin{minipage}{10cm}
			\tableofcontents[currentsection]
		\end{minipage}
	}
}

% TOC without subsection
\setcounter{tocdepth}{1} % only-- part,chapters,sections 

% bullet points: rectangles
\useinnertheme{rectangles}
\setbeamercolor{itemize item}{fg=RUBLightGreen}
\setbeamercolor{itemize subitem}{fg=RUBLightGreen}
% enumerate: blue background for better readability
\setbeamercolor{item projected}{bg=RUBDarkBlue}

% make boxes (example, block, etc.) background lighter for readability
\setbeamercolor{block title}{%
	use=normal text,
	fg=normal text.fg,
	bg=normal text.bg!90!fg % lighter background in block title
}
\setbeamercolor{block body}{
	use={block title, normal text},
	bg=block title.bg!30!normal text.bg % lighter background in block body
}


% RUB colors in blocks
\setbeamercolor{block title alerted}{%
	use={block title, alerted text},
	bg=RUBDarkBlue,
	%fg=RUBLightYellow % looks bad
	fg=white % better contrast
}

\setbeamercolor{block title example}{%
	use={block title, example text},
	fg=RUBLightGreen
}


% ------- end of RUB specifics ----------

% all itemize with pause by default
%\beamerdefaultoverlayspecification{<+->}


% typeset mathematics on serif
\usefonttheme[onlymath]{serif}

% better bibliography using biber as backend
\usepackage[natbib=true,backend=biber,style=authoryear-icomp,maxbibnames=30,maxcitenames=9,uniquelist=false,giveninits=true,doi=false,url=false,dashed=false,isbn=false]{biblatex}
% shared bibliography
\addbibresource{../bibliography.bib}
% disable "ibid" for repeated citations
\boolfalse{citetracker}



\usepackage{xspace}


% for derivatives, https://tex.stackexchange.com/a/412442
\usepackage{physics}

\usepackage{tikz}
\usetikzlibrary{matrix, positioning}
\usetikzlibrary{angles,quotes} % for angles
\usetikzlibrary{backgrounds} % background
\usetikzlibrary{decorations.pathreplacing} % curly braces
\usetikzlibrary{calligraphy}
\usetikzlibrary{calc} % for neural nets

% for plotting functions
\usepackage{pgfplots}
\usepgfplotslibrary{dateplot}

% sub-figures
\usepackage{caption}
\usepackage{subcaption}

% book tabs
\usepackage{booktabs}


% argmin, argmax
\usepackage{amsmath}
\DeclareMathOperator*{\argmax}{arg\!\max}
\DeclareMathOperator*{\argmin}{arg\!\min}
% softmax
\DeclareMathOperator*{\softmax}{soft\!\max}
% Mask
\DeclareMathOperator*{\mask}{mask}

% bold math
\usepackage{bm}

% for \mathclap
\usepackage{mathtools}

% algorithms
\usepackage[noend]{algpseudocode}


% for neurons and layers in tikz
\tikzset{
	neuron/.style={draw, rectangle, inner sep=2pt, minimum width=0.75cm, fill=blue!20},
	param/.style={draw, rectangle, inner sep=2pt, minimum width=0.75cm, fill=green!20},
	constant/.style={draw, rectangle, inner sep=2pt, minimum width=0.75cm, fill=black!15},
	% for citation nodes right top
	ref/.style={anchor = north east, text width=7.8cm, yshift=-1.3cm, xshift=-0.2cm, scale=0.5},
	state/.style={rectangle, inner sep=2pt, minimum width=0.75cm, fill=black!5},
}

% added in lecture 10
\tikzset{
	mtx/.style={
		matrix of math nodes,
		left delimiter={[}, right delimiter={]}
	},
	hlt/.style={opacity=0.1, line width=4 mm, line cap=round},
	hltr/.style={opacity=0.5, rounded corners=2pt, inner sep=-1pt}
}

% for strike-through text (added in Lecture 06)
\usepackage[normalem]{ulem}

% added in Lecture 7
% RNN
\DeclareMathOperator*{\rnn}{RNN}
% RNN star
\DeclareMathOperator*{\rnnstar}{RNN^{*}}
% bi-RNN
\DeclareMathOperator*{\birnn}{biRNN}


% added in Lecture 9
\usetikzlibrary{fit} % for hightligting by calling "fit"

% algorithms
\usepackage[noend]{algpseudocode}


\title{Privacy-Preserving Natural Language Processing}
\subtitle{Lecture 10 : Real-World DP Application and the Epsilon Registry}
\date{July 17, 2025}
\author{Dr.\ Erion {\c C}ano}
\institute{
\texttt{www.trusthlt.org} \\
Chair of Trustworthy Human Language Technologies (TrustHLT) \\
Ruhr University Bochum \& Research Center Trustworthy Data Science and Security}

\begin{document}


\maketitle


\section{Tech Giants Applying DP}

%-----------------------------------------------------------------------

\begin{frame}{Apple User Data from macOS and iOS}

Tons of user data being collected by Apple:

\begin{itemize} \setlength\itemsep{2mm}
%
\item Typing information for QuickType suggestions
%
\item Emojis for emoji suggestions
%
\item Web navigation data from Safari
%
\end{itemize}

Apple uses local DP with different $\varepsilon$ values to protect user privacy in those data. \vskip1mm

Full details are reported {\it \href{https://docs-assets.developer.apple.com/ml-research/papers/learning-with-privacy-at-scale.pdf}{here}}.

\end{frame}

%-----------------------------------------------------------------------

\begin{frame}{Meta and Facebook User Interactions}

Dataset of user interactions with web pages shared on FB is publicly available. \vskip1mm

Examples: \vskip1mm 

\begin{itemize} \setlength\itemsep{2mm}
%
\item John shared URL foo
%
\item Mary views post with URL bar
%
\end{itemize}

Meta uses $(\varepsilon,\delta)$ DP to protect each interaction.

\begin{tikzpicture}[overlay, remember picture]
\node at (current page.north east)[ref] {
\fullcite{DVN/TDOAPG_2020} \par};
\end{tikzpicture}

\end{frame}

%-----------------------------------------------------------------------

\begin{frame}{Microsoft and Windows Telemetry}

User data collections starting (massively) with Windows 10. \vskip1mm

Examples: \vskip1mm 

\begin{itemize} \setlength\itemsep{2mm}
%
\item Application crash reports. 
%
\item User time spent on certain apps.
%
\item Typing, location, connected devices, etc. 
%
\end{itemize}

Microsoft uses local DP with $\varepsilon = 1.67$ to protect it before using / sharing / selling.

\begin{tikzpicture}[overlay, remember picture]
\node at (current page.north east)[ref] {
\fullcite{ding2017collecting} \par};
\end{tikzpicture}

\end{frame}

%-----------------------------------------------------------------------

\begin{frame}{Google Shopping and Trends}

Product page view counts for prioritizing page crawing. 

\begin{itemize} \setlength\itemsep{2mm}
%
\item Collected from users on daily basis. 
%
\item Streammed using DP-SQLP
%
\item Protected using $(\varepsilon, \delta)$ DP with $\varepsilon = 1$ and $\delta = 10^{-9}$
%
\end{itemize}

Search trends used to show related queries. 
 
\begin{itemize} \setlength\itemsep{2mm}
%
\item Streammed using DP-SQLP
%
\item Protected using $(\varepsilon, \delta)$ DP with $\varepsilon = 2$ and $\delta = 10^{-10}$
%
\end{itemize}

\end{frame}

%-----------------------------------------------------------------------

\begin{frame}{Linkedin Audience Engagement}

Interactive query system that allows marketers to get information about LinkedIn users engaging with their content. 

\begin{itemize} \setlength\itemsep{2mm}
%
\item Each analyst sends queries to the system 
%
\item Each query returns $(\varepsilon, \delta)$ DP with $\varepsilon = 0.15$ and $\delta = 10^{-10}$
%
\item Monthly limit (budget) for queries is $\varepsilon = 34.9$ and $\delta = 7 \cdot 10^{-9}$
%
\item Extra measures to prevent averaging attacks 
\end{itemize}

\begin{tikzpicture}[overlay, remember picture]
\node at (current page.north east)[ref] {
\fullcite{DBLP:journals/corr/abs-2002-05839} \par};
\end{tikzpicture}

\end{frame}

%-----------------------------------------------------------------------

\begin{frame}{Labor Market Insights}

Data from LinkedIn to measure trends in people changing their occupation. \vskip2mm 

\begin{itemize} \setlength\itemsep{2mm}
%
\item Companies who are hiring, using with $\varepsilon = 14.4$ and $\delta = 1.2 \cdot 10^{-9}$
%
\item Available jobs, protecting hiring event with $\varepsilon = 14.4$ and $\delta = 1.2 \cdot 10^{-9}$
%
\item Required skills, protecting each user's skills information with $\varepsilon = 0.3$ and $\delta = 3 \cdot 10^{-10}$
%
\item Extra measures to prevent averaging attacks 
\end{itemize}

\begin{tikzpicture}[overlay, remember picture]
\node at (current page.north east)[ref] {
\fullcite{rogers2020membersapproachenablinglinkedins} \par};
\end{tikzpicture}

\end{frame}

%-----------------------------------------------------------------------

\section{Other Organizations Applying DP}

%-----------------------------------------------------------------------

\begin{frame}{USA Census Bureau}

\textbf{County Business Patterns} \vskip1mm

\begin{itemize} \setlength\itemsep{2mm}
%
\item Business establishments in the USA
%
\item Different types of data (including finances) 
%
\item Protected with various $(\varepsilon, \delta)$ parameters based on:
%
	\begin{itemize} \setlength\itemsep{2mm}
	%
	\item Financial assets
	%
	\item Annual payroll 
	%
	\item Number of employs 
	%
	\end{itemize}
%
\end{itemize}

Full details are reported {\it \href{https://www.census.gov/topics/business-economy/disclosure/data/tables/cbp-privacy-demonstration-tables.html}{here}}.

\end{frame}

%-----------------------------------------------------------------------

\begin{frame}{USA Census Bureau}

\textbf{The 2020 USA Census} \vskip1mm

\begin{itemize}
%
\item Collecting demographic information about USA population \vskip1mm
%
\item Datasets protected with various $(\varepsilon, \delta)$ parameters \vskip1mm
%
	\begin{itemize} \setlength\itemsep{2mm}
	%
	\item {\it \href{https://www.census.gov/programs-surveys/decennial-census/about/rdo/summary-files.html}{Redistricting}} data protected with $\varepsilon = 13.64$
	%
	\item {\it \href{https://www.census.gov/data/tables/2023/dec/2020-census-dhc.html}{Demographic Housing}} protected with $\varepsilon = 19.46$
	%
	\item Racial and ethnical categories protected with $\varepsilon = 45.68$
	%
	\item Suplemental household statistics protected with $\varepsilon = 12.74$
	%
	\end{itemize}
%
\end{itemize}

\end{frame}

%-----------------------------------------------------------------------

\begin{frame}{USA Census Bureau}

\textbf{Post-Secondary Employment Outcomes} \vskip2mm

\begin{itemize} \setlength\itemsep{2mm}
%
\item Earnings and employment of college graduates
%
\item Categorized by degree level, degree major, and post-secondary institution
%
\item Matching university granscript data with national database of jobs
%
\item Each person protected with $\varepsilon = 1.5$
%
\end{itemize}

Full details are reported {\it \href{https://lehd.ces.census.gov/doc/PSEOTechnicalDocumentation.pdf}{here}}.

\end{frame}

%-----------------------------------------------------------------------

\begin{frame}{Wikimedia Foundation}

\textbf{Page view statistics} \vskip2mm

\begin{itemize} \setlength\itemsep{2mm}
%
\item Distinct users visiting Wiki pages each day from each country
%
\item July 1, 2015 to Feb. 8, 2017 protected with $\varepsilon = 1$ on 300 pages view per day
%
\item Feb. 9, 2017 to Feb. 5, 2023 protected with $\varepsilon = 1$ on 30 pages view per day
%
\item Feb. 6, 2023 onwards protected with $\varepsilon = 0.72$ and $\delta = 10^{-5}$
%
\end{itemize}

\begin{tikzpicture}[overlay, remember picture]
\node at (current page.north east)[ref] {
\fullcite{adeleye2023publishingwikipediausagedata} \par};
\end{tikzpicture}

\end{frame}

%-----------------------------------------------------------------------

\begin{frame}{Wikimedia Foundation}

\textbf{Editor statistics} \vskip2mm

\begin{itemize} \setlength\itemsep{2mm}
%
\item Statistics about editor activity by project and country \vskip2mm
%
\item Monthly published data protected with $\varepsilon = 2$ on editor-project-country-month
%
\item Weekly published data protected with $\varepsilon = 2$ on editor-project-country-week
%
\item One-off release for Russian editors, protected with $\varepsilon = 0.1$ 
%
\end{itemize}

\end{frame}

%-----------------------------------------------------------------------

\begin{frame}{Other Applications}

\begin{itemize} \setlength\itemsep{2mm}
%
\item Brave uses distributed DP for collecting user analytics
%
\item Spectus published a dashboard with mobility trends during Hurricane Irma
%
\item Microsoft Assistive AI collects user data about Office Tools and suggests automatic replies
%
\item Google uses DP and federated learning for collecting data and training models for improved text selection and copying on Android
%
\end{itemize}

\end{frame}

%-----------------------------------------------------------------------

\section{Standardizations and the Epsilon Registry}

%-----------------------------------------------------------------------

\begin{frame}{Background: Evolution of Information Security}

\textbf{Early days, until the '60s} \vskip2mm

\begin{itemize} \setlength\itemsep{2mm}
%
\item Mostly physical security: focus on access to mainframes 
%
\item Limited threats: computers were rare and isolated
%
\item Computers were mostly used for scientific computations; little or no personal data involved
%
\item Basic enkryption: simple ciphers like Caesar sipher were used to ensure message confidentiality
%
\end{itemize}

\end{frame}

%-----------------------------------------------------------------------

\begin{frame}{Background: Evolution of Information Security}

\textbf{Rise of the networks, '60s - '90s} \vskip2mm

\begin{itemize} \setlength\itemsep{1mm}
%
\item Computer networks and remote access become common
%
\item Unauthorized network / system access and data breaches become popular concerns
%
\item Cybersecurity emerges as an important discipline
%
\item Many remote access protocols with encryption introduced
% 
\item Popular symetric cyphers of 64 bits like DES 
%
\item Public key cryptography comes out
%
\end{itemize}

\end{frame}

%-----------------------------------------------------------------------

\begin{frame}{Background: Evolution of Information Security}

\textbf{The internert era, '00s - present} \vskip2mm

\begin{itemize} \setlength\itemsep{2mm}
%
\item Widespread internet usage, online shopping, personal data sharing
%
\item Surge in cybercrime, viruses, worms, ransomware, phishing, DDOS attacks
%
\item Evolution of security measures: VPNs, firewalls, antiviruses, intrusion detection systems
%
\item Evolution of cryptography: 128 bit cyphers like AES become mandatory
%
\end{itemize}

\end{frame}

%-----------------------------------------------------------------------

\begin{frame}{Information Security Standards}

\begin{itemize} \setlength\itemsep{4mm}
%
\item \textbf{ISO/IEC 27001}: Standards for establishing, implementing, maintaining and improving an ISMS
%
\item \textbf{ISO/IEC 27002}: provides a code of practice for information security controls, offering guidance for implementing security policies
%
\item \textbf{ISO/IEC 27005}: focuses on information security risk management, helping organizations in assessing and managing risks
%
\end{itemize}

\end{frame}

%-----------------------------------------------------------------------

\begin{frame}{Information Privacy vs. Security}

\begin{itemize} \setlength\itemsep{3mm}
%
\item Privacy is inherently a more complex concept: who, what, when, where, why
%
\item Private information is highly context dependent which makes it difficult to protect
%
\item Lack of motivation (unless enforced by law) to protect privacy because of utility (profitability) loss
%
\item Can privacy evolution and standardization follow the same path as security?
%
\end{itemize}

\end{frame}

%-----------------------------------------------------------------------

\begin{frame}{The Epsilon Registry}

\begin{itemize} \setlength\itemsep{3mm}
%
\item DP is being successfully implemented in industry, public sector, academia, etc. 
%
\item Still, little understanding of the optimal $\varepsilon$ values for certain systems, purposes, data types, etc.
%
\item No clear consensus among practicioneers about:
%
	\begin{itemize} \setlength\itemsep{2mm}
	%
	\item How to approach the key implementation decisions
	%
	\item How to decide about the right privacy guarantees
	%
	\item How to choose $\varepsilon$ levels 
	%
	\end{itemize}
%
\end{itemize}

\begin{tikzpicture}[overlay, remember picture]
\node at (current page.north east)[ref] {
\fullcite{dwork_differential_2019} \par};
\end{tikzpicture}

\end{frame}

%-----------------------------------------------------------------------

\begin{frame}{The Epsilon Registry}

Could a public Epsilon Registry help?

\begin{itemize} \setlength\itemsep{3mm}
%
\item Communal body of knowledge about DP implementations
%
\item Used by stakeholders, public servants, academics, etc. 
%
\item Serve as guideline for identifying and adopting judicious DP implementations
%
\end{itemize}

\begin{tikzpicture}[overlay, remember picture]
\node at (current page.north east)[ref] {
\fullcite{dwork_differential_2019} \par};
\end{tikzpicture}

\end{frame}

%-----------------------------------------------------------------------

\begin{frame}{Questions...?}

\pause 

\centering
\textbf{THANK YOU...!}

\end{frame}

%------------------------------------------------------------------------

\end{document}

